\section{Arquitetura do Sistema}

A arquitetura do NSS utiliza quatro componentes executados separadamente: frontend, backend Java/Spring Boot, pipeline Python e PostgreSQL.

\placeholderfigure{arquitetura-alto-nivel}{7cm}{Arquitetura de alto nível do Núcleo de Situação de Saúde.}

\subsection{Fluxo de consulta}

O fluxo normal de uma consulta é:

\begin{enumerate}
    \item o usuário interage com o frontend;
    \item o frontend envia uma requisição HTTP para a API Java;
    \item o Java autentica e autoriza a operação, quando aplicável;
    \item a camada de serviço aplica as regras de negócio;
    \item o backend consulta o PostgreSQL;
    \item o resultado é transformado em DTO e retornado ao frontend.
\end{enumerate}

\subsection{Fluxo de atualização de dados}

A atualização dos dados ocorre separadamente do fluxo de consulta:

\begin{enumerate}
    \item o pipeline Python obtém dados do SINAN por meio do PySUS;
    \item os dados brutos são preservados em Bronze;
    \item transformações e padronizações produzem Silver;
    \item agregações necessárias à aplicação produzem Gold;
    \item os resultados são validados;
    \item os dados consolidados são persistidos no PostgreSQL.
\end{enumerate}

\begin{regra}[Pipeline fora do caminho crítico de consulta]
Uma consulta do usuário não deve depender de uma nova coleta ou processamento do SINAN. A aplicação consulta dados já persistidos e validados no PostgreSQL.
\end{regra}

\subsection{Responsabilidades}

\begin{tabularx}{\textwidth}{@{}p{3.2cm}X@{}}
\toprule
\textbf{Componente} & \textbf{Responsabilidade} \\
\midrule
Frontend & Interface, visualizações, mapas, filtros e interação com o usuário. \\
Java/Spring Boot & API REST, autenticação, autorização, regras de negócio e acesso ao banco. \\
Python/Pipeline & Ingestão, transformação, agregação, validação e atualização dos dados. \\
PostgreSQL & Persistência compartilhada dos dados consolidados e das informações da aplicação. \\
\bottomrule
\end{tabularx}
